\section{Appendix}

The complete source code is organized below by architectural layer. All source files are available in the project GitHub repository.

\textbf{GitHub Repository:} \url{https://github.com/mcittkmims/es-labs}

% ─────────────────────────────────────────────────────────────────────────────
\subsection{Application Layer}

\subsubsection{Main Entry Point}

\begin{lstlisting}[language=C++, caption=\texttt{src/main.cpp} --- lab selector (Lab 4.1 excerpt), label=lst:appendix_main]
// Automatically select the active lab based on build flags
#elif defined(LAB4_1)
    #include "lab4_1_main.h"

void setup() {
#elif defined(LAB4_1)
    lab4_1Setup();
}
void loop() {
#elif defined(LAB4_1)
    lab4_1Loop();
}
\end{lstlisting}

\subsubsection{Lab 4.1 Entry Point}

\begin{lstlisting}[language=C++, caption=\texttt{lab4\_1/lab4\_1\_main.h}, label=lst:appendix_lab_h]
void lab4_1Setup();   // Initialize hardware, shared state, and FreeRTOS tasks
void lab4_1Loop();    // Empty - all logic in FreeRTOS tasks
\end{lstlisting}

\begin{lstlisting}[language=C++, caption=\texttt{lab4\_1/lab4\_1\_main.cpp}, label=lst:appendix_lab_cpp]
#include "lab4_1_main.h"
#include "actuator_data.h"
#include "task_input.h"
#include "task_actuator.h"
#include "task_display.h"
#include <Arduino_FreeRTOS.h>
#include "StdioSerial.h"

void lab4_1Setup() {
    stdioSerialInit(9600);
    printf("Lab 4.1 -- Binary Actuator Control\r\n");
    printf("Debounce: %u ms, Persist: %u confirms\r\n",
           DEBOUNCE_PERIOD_MS, PERSIST_COUNT);

    actuatorDataInit();

    xTaskCreate(vTaskInputReader, "Input",   TASK_INPUT_STACK,
                NULL, TASK_INPUT_PRIORITY,    NULL);
    xTaskCreate(vTaskActuator,    "Actuator",TASK_ACTUATOR_STACK,
                NULL, TASK_ACTUATOR_PRIORITY, NULL);
    xTaskCreate(vTaskDisplay,     "Display", TASK_DISPLAY_STACK,
                NULL, TASK_DISPLAY_PRIORITY,  NULL);
}

void lab4_1Loop() { /* All logic in FreeRTOS tasks */ }
\end{lstlisting}

\subsubsection{Shared State and Constants}

\begin{lstlisting}[language=C++, caption=\texttt{lab4\_1/actuator\_data.h}, label=lst:appendix_data_h]
#define TASK_INPUT_PERIOD_MS      10U
#define TASK_ACTUATOR_PERIOD_MS   75U
#define TASK_DISPLAY_PERIOD_MS   500U
#define STDIO_REPORT_EVERY_N       4U
#define PIN_ACTUATOR_LED           7U
#define DEBOUNCE_PERIOD_MS       100U
#define PERSIST_COUNT              3U

typedef struct {
    ActuatorControl::State pendingCmd;
    ActuatorControl::State actuatorState;
    bool                   isDebouncing;
    uint8_t                confirmCount;
    uint32_t               toggleCount;
    SemaphoreHandle_t      mutex;
} ActuatorShared;

extern ActuatorShared g_actuator;
void actuatorDataInit();
\end{lstlisting}

\subsubsection{Input Reader Task}

\begin{lstlisting}[language=C++, caption=\texttt{lab4\_1/task\_input.cpp}, label=lst:appendix_input]
void vTaskInputReader(void *params) {
    char buf[32]; uint8_t len = 0;
    printf("Commands: ON | OFF\r\n> ");
    TickType_t lastWake = xTaskGetTickCount();
    while (true) {
        while (Serial.available()) {
            char c = (char)Serial.read();
            Serial.write(c);
            if (c == '\r' || c == '\n') {
                if (len > 0) {
                    buf[len] = '\0'; normalizeInput(buf);
                    if (strcmp(buf, "ON") == 0) {
                        // take mutex, write pendingCmd = ON, give mutex
                    } else if (strcmp(buf, "OFF") == 0) {
                        // take mutex, write pendingCmd = OFF, give mutex
                    } else {
                        printf("  [ERR] Unknown: '%s'\r\n", buf);
                    }
                    len = 0; printf("> ");
                }
            } else if (len < sizeof(buf)-1) buf[len++] = c;
        }
        vTaskDelayUntil(&lastWake, pdMS_TO_TICKS(TASK_INPUT_PERIOD_MS));
    }
}
\end{lstlisting}

\subsubsection{Actuator Control Task}

\begin{lstlisting}[language=C++, caption=\texttt{lab4\_1/task\_actuator.cpp}, label=lst:appendix_actuator]
static ActuatorControl s_actuator(PIN_ACTUATOR_LED,
                                   DEBOUNCE_PERIOD_MS, PERSIST_COUNT);

void vTaskActuator(void *params) {
    s_actuator.init();
    TickType_t lastWake = xTaskGetTickCount();
    while (true) {
        // 1. Read raw command
        ActuatorControl::State rawCmd;
        xSemaphoreTake(g_actuator.mutex, pdMS_TO_TICKS(5));
        rawCmd = g_actuator.pendingCmd;
        xSemaphoreGive(g_actuator.mutex);

        // 2. Run signal conditioning pipeline + drive GPIO
        s_actuator.setRawCommand(rawCmd);
        s_actuator.update();

        // 3. Write conditioned state back to shared memory
        xSemaphoreTake(g_actuator.mutex, pdMS_TO_TICKS(5));
        if (s_actuator.getState() != g_actuator.actuatorState)
            g_actuator.toggleCount++;
        g_actuator.actuatorState = s_actuator.getState();
        g_actuator.isDebouncing  = s_actuator.isDebouncing();
        g_actuator.confirmCount  = s_actuator.getConfirmCount();
        xSemaphoreGive(g_actuator.mutex);

        vTaskDelayUntil(&lastWake, pdMS_TO_TICKS(TASK_ACTUATOR_PERIOD_MS));
    }
}
\end{lstlisting}

\subsubsection{Display and Reporting Task}

\begin{lstlisting}[language=C++, caption=\texttt{lab4\_1/task\_display.cpp}, label=lst:appendix_display]
static LcdDisplay s_lcd(0x27, 16, 2);

void vTaskDisplay(void *params) {
    s_lcd.init();
    s_lcd.showTwoLines("Actuator Ctrl", "Initializing...");
    uint8_t tick = 0;
    TickType_t lastWake = xTaskGetTickCount();
    while (true) {
        // Read shared state, update LCD, print STDIO report every 4 ticks
        // LCD line 1: "Actuator: ON " or "Actuator: OFF"
        // LCD line 2: "STABLE    " or "DEBOUNCING"
        vTaskDelayUntil(&lastWake, pdMS_TO_TICKS(TASK_DISPLAY_PERIOD_MS));
    }
}
\end{lstlisting}

% ─────────────────────────────────────────────────────────────────────────────
\subsection{Service Layer}

\subsubsection{ActuatorControl}

\begin{lstlisting}[language=C++, caption=\texttt{lib/ActuatorControl/ActuatorControl.h}, label=lst:appendix_actuatorctrl_h]
class ActuatorControl {
public:
    enum class State { OFF, ON };
    ActuatorControl(uint8_t pin, uint32_t debouncePeriodMs=100,
                    uint8_t persistCount=3);
    void  init();
    void  setRawCommand(State cmd);
    void  update();                   // run pipeline + drive GPIO
    State getState()      const;      // conditioned hardware state
    State getRawCommand() const;
    bool  isDebouncing()  const;
    uint8_t getConfirmCount()         const;
    uint8_t getRequiredConfirmCount() const;
private:
    Led               _led;
    BinaryConditioner _conditioner;
    State             _rawCommand;
    State             _currentState;
};
\end{lstlisting}

\begin{lstlisting}[language=C++, caption=\texttt{lib/ActuatorControl/ActuatorControl.cpp} (key method), label=lst:appendix_actuatorctrl_cpp]
void ActuatorControl::update() {
    bool rawBool     = (_rawCommand == State::ON);
    bool conditioned = _conditioner.process(rawBool);
    State newState   = conditioned ? State::ON : State::OFF;
    if (newState != _currentState) {
        _currentState = newState;
        if (_currentState == State::ON) _led.turnOn();
        else                            _led.turnOff();
    }
}
\end{lstlisting}

% ─────────────────────────────────────────────────────────────────────────────
\subsection{ECAL Layer}

\subsubsection{BinaryConditioner}

\begin{lstlisting}[language=C++, caption=\texttt{lib/BinaryConditioner/BinaryConditioner.h}, label=lst:appendix_conditioner_h]
class BinaryConditioner {
public:
    BinaryConditioner(uint32_t debouncePeriodMs=100, uint8_t persistCount=3);
    bool    process(bool rawInput);   // run 3-stage pipeline
    bool    getConditionedState() const;
    bool    getRawState()         const;
    bool    getPendingState()     const;
    bool    isDebouncing()        const;
    uint8_t getConfirmCount()     const;
    uint8_t getRequiredConfirmCount() const;
    void    reset();
private:
    uint32_t _debouncePeriodMs;
    uint8_t  _persistCount;
    bool     _rawState, _pendingState, _conditionedState;
    uint32_t _debounceStartMs;
    bool     _inDebounce;
    uint8_t  _confirmCount;
};
\end{lstlisting}

\begin{lstlisting}[language=C++, caption=\texttt{lib/BinaryConditioner/BinaryConditioner.cpp}, label=lst:appendix_conditioner_cpp]
bool BinaryConditioner::process(bool rawInput) {
    bool saturated = rawInput ? true : false;  // Stage 1: Saturation
    _rawState = saturated;

    if (saturated != _pendingState) {
        _pendingState    = saturated;           // restart on change
        _debounceStartMs = millis();
        _confirmCount    = 1;
    } else if (_pendingState != _conditionedState) {
        _confirmCount++;                        // Stage 3: accumulate
    }

    if (_pendingState != _conditionedState) {
        _inDebounce = true;
        uint32_t elapsed = millis() - _debounceStartMs;
        // Stage 2+3: commit only when BOTH satisfied
        if (_confirmCount >= _persistCount && elapsed >= _debouncePeriodMs) {
            _conditionedState = _pendingState;
            _inDebounce = false; _confirmCount = 0;
        }
    } else { _inDebounce = false; }
    return _conditionedState;
}
\end{lstlisting}

% ─────────────────────────────────────────────────────────────────────────────
\subsection{MCAL Layer}

\subsubsection{Led Driver}

\begin{lstlisting}[language=C++, caption=\texttt{lib/Led/Led.h} (interface summary), label=lst:appendix_led_h]
class Led {
public:
    Led(uint8_t pin);
    void init();       // pinMode OUTPUT, initial OFF
    void turnOn();     // digitalWrite HIGH
    void turnOff();    // digitalWrite LOW
    void toggle();
    bool isOn() const;
private:
    uint8_t ledPin;
    bool    state;
};
\end{lstlisting}

% ─────────────────────────────────────────────────────────────────────────────
\subsection{PlatformIO Configuration}

\begin{lstlisting}[caption=\texttt{platformio.ini} --- Lab 4.1 environment, label=lst:appendix_platformio]
[env:lab4_1]
platform = atmelavr
board = megaatmega2560
framework = arduino
monitor_speed = 9600
build_src_filter = +<*> +<../lab/lab4_1/*>
build_flags = -I lab/lab4_1 -DLAB4_1
lib_ldf_mode = deep+
lib_deps =
    feilipu/FreeRTOS
    marcoschwartz/LiquidCrystal_I2C@^1.1.4
\end{lstlisting}
